\scanpage{337}
\mspage{76}
% archive sheet number: 336
% ms. 1–4
\unsure{correspondante|correspondant}, \deleted{\unsure{on}} on trouve une homotopie de $a$ \unsure{à|et} $b$ de type $\tau_{\mathrm{red}}$, dont la longueur est égale à la ``\uline{longueur utile}'' du type $\tau$ de $c$ (ce qui justifie la terminologie).
% ms. 5
Notons que si
% ms. 6
\[\tau=\Delta_{\varepsilon\deleted{\unsure{+}}}^{\ell_1}\,\Delta_{-\varepsilon\deleted{\unsure{2}}}^{\ell_2}\,\Delta_\varepsilon^{\ell_3}\cdots\Delta_{(-1)^{\unsure{r|n}}\varepsilon}^{\ell_{\deleted{\unsure{n}}r}},\]
% ms. 7
alors
% ms. 8
\[\tau_{\mathrm{red}}=\underbrace{\Delta_\varepsilon\,\Delta_{-\varepsilon}\,\Delta_{+\varepsilon}\cdots\Delta_{(-1)^r\varepsilon}}_{r},\]
% ms. 9–10
où $r$ est la longueur utile de $\tau$ (la longueur \unsure{pas} utile étant $n=\sum_{i\in[1,r]}\ell_i$).

% ms. 11–12
Quand on par\deleted{\unsure{t}}le d'une homotopie entre \uline{foncteurs}
% ms. 12
\[X\longrightarrow Y,\]
% ms. 13–16
il est entendu qu'il s'agit d'homotopies dans la catégorie $\underline{\mathrm{Hom}}(X\unsure{,|;}Y)$\unsure{.|;} (Ou, suivant le contexte, dans \deleted{\unsure{telle}}\unsure{une} sous-catégorie précisée de $\underline{\mathrm{Hom}}(X,Y)$.)
% ms. below 16
\figure{\begin{tikzpicture}[x=1cm,y=1cm]\draw[line width=.55pt] (0,0) .. controls (.18,-.26) and (.62,-.67) .. (.66,-.91);\end{tikzpicture}}
